\documentclass[11pt,a4paper]{article}

% ======================================================================
%  Operational Coherence Maintenance: closure note — Version 2 (July 2026)
%  v1: December 2025 (ai.viXra:2512.0071v1).
%  Changes: the imported theorem is re-stated per the corrected 0061 v2
%  (v1 imported the vacuous unlinked-pairs form); the rate-inheritance
%  hinge is updated from "open" to "partially resolved" (proved in the
%  quasi-free local-sink class, 0064 v2; the omega=0-floor scenario
%  realized within the Davies class, 0070 v2); proxy-validation lesson
%  added; references updated to the v2 series and companion repository.
% ======================================================================

\usepackage[utf8]{inputenc}
\usepackage[T1]{fontenc}
\usepackage{amsmath,amssymb,amsthm,mathtools}
\usepackage[margin=1.05in]{geometry}
\usepackage[colorlinks=true,linkcolor=blue,citecolor=blue,urlcolor=blue]{hyperref}

\theoremstyle{plain}
\newtheorem{theorem}{Theorem}[section]
\theoremstyle{definition}
\newtheorem{definition}[theorem]{Definition}
\newtheorem{hypothesis}[theorem]{Hypothesis}
\theoremstyle{remark}
\newtheorem{remark}[theorem]{Remark}

\newcommand{\Tr}{\operatorname{Tr}}
\newcommand{\kB}{k_{\mathrm{B}}}

\title{Operational Coherence Maintenance:\\
Proven Results, Conditional Interfaces, and Open Dynamical Gaps}
\author{Lluis Eriksson\\ \small Independent Researcher\\
\small \texttt{lluiseriksson@gmail.com}}
\date{July 4, 2026 \quad (v2; v1: December 2025) \quad ai.viXra:2512.0071}

\begin{document}

\maketitle

\begin{center}
\fbox{\parbox{0.92\textwidth}{\small
\textbf{Editorial status: companion note.} Architectural/closure overview of
the coherence-maintenance series; no new technical results. The series
consists of ai.viXra:2512.0060, 0061, 0064, 0070, and the RIP framing note
2512.0072, all at v2 (2026), each with verification material; papers, scripts,
and reference logs are distributed in a single companion repository. This note
is the map.}}
\end{center}

\begin{abstract}
Maintaining quantum coherence against uncontrolled open-system dynamics is an
operational control task with unavoidable thermodynamic cost. In finite
dimensions, explicit lower bounds on the minimal power required to stabilize
coherence can be derived under standard Markovian assumptions, independently
of geometric or field-theoretic structure. At the same time, static
correlations in gapped systems are geometrically suppressed, raising the
question of how such suppression influences dynamical decoherence rates and,
consequently, coherence-maintenance power. Bridging the two domains requires
dynamical input that static clustering alone does not provide.

This note introduces no new technical results. It provides a logical closure
of the program by separating (i) results proven without additional structure,
(ii) conditional interfaces, and (iii) dynamical hypotheses --- and, new in
this version, it \emph{updates the status of the central hinge}. In v1, rate
inheritance --- the relation between static correlation envelopes and
effective decoherence rates --- was identified as the unique unresolved
hinge. Since then it has been \emph{partially resolved} in both directions
anticipated by v1's scenario analysis: it is now a derived, frequency-resolved
law in an exactly solvable quasi-free local-sink class (2512.0064 v2), and the
failure scenario through near-zero-frequency channels has been realized within
the Davies model class, with the persistent floor computed and the secular
nonlocality of that construction quantified (2512.0070 v2). The imported
maintenance bound is restated in its corrected v2 form (2512.0061 v2), whose
v1 formulation was vacuous. The framework's design goal --- robustness under
partial refutation --- has thus been exercised in practice, twice.
\end{abstract}

\section{Introduction: why a closure note is useful}

The loss of quantum coherence under uncontrolled dynamics is a central feature
of open quantum systems. Decoherence theory explains passive suppression of
interference at the level of reduced dynamics, but is largely silent on a
complementary operational question: what control resources are required to
\emph{prevent} this loss and maintain a target coherent state in steady time?

Coherence maintenance can be formulated as a thermodynamic control task. Under
explicit control models --- battery-assisted thermal operations --- one derives
lower bounds on the minimal power required to counteract uncontrolled dynamics
\cite{MaintWork}. These bounds depend on dynamical quantities such as the
intrinsic coherence-loss rate, and do not rely on spatial structure, field
theory, or geometric separation. Independently, static correlations in gapped
systems are suppressed with distance, often exponentially
\cite{HastingsKoma}, with Bessel-type envelopes in continuum massive regimes
\cite{Clustering}. It is natural to ask whether geometric suppression of
static correlations constrains dynamical decoherence rates and the power
required to maintain coherence across an operational interface.

Five manuscripts address different parts of this question: algebraic
kinematics and collar-suppressed correlations \cite{Clustering}, operational
maintenance-power bounds \cite{MaintWork}, exact dynamical rate inheritance
and its proxy pathologies \cite{HCut}, a Davies-class stress test with
ceiling/floor envelopes \cite{Stress}, and the formal statement of the
rate-inheritance principle itself, with its proxy-validated falsification
protocol \cite{RIP}. The purpose of this note is to make the
logical structure explicit --- what is proved, what is conditional, what
remains open --- and, in this version, to record how the open hinge of v1 has
since moved.

\begin{remark}[What changed since v1, in one paragraph]\label{rem:whatmoved}
Every paper in the series was revised in 2026 after an exact-verification
campaign, and this note inherits those corrections: the reconstruction theorem
of \cite{Clustering} is now stated for conditional reattachment (not the Petz
map, which fails marginal preservation for correlated references); the work
theorem of \cite{MaintWork} corrected a sign error and replaced a vacuous
extra-power definition; the dynamical evidence of \cite{HCut} v1 was withdrawn
and replaced by an exact law; and the envelope of \cite{Stress} v1 was split
into a ceiling and the floor that a no-go actually needs. The architecture of
v1 of this note survived all four revisions unchanged --- which is what it was
for --- but its statements of the imported results and of the status of the
hinge did not, and are updated below.
\end{remark}

\section{Layer I: proved core statements (finite-dimensional,
model-independent)}

\begin{definition}[Relative entropy, energy pinching, coherence-loss rate]
For density operators $\rho, \sigma$ with $\mathrm{supp}(\rho) \subseteq
\mathrm{supp}(\sigma)$, $S(\rho\|\sigma) := \Tr[\rho(\log\rho - \log\sigma)]$
($+\infty$ otherwise). With $H = \sum_n E_n\Pi_n$, the energy pinching is
$\Delta[\rho] := \sum_n \Pi_n\rho\Pi_n$, the coherence is $C(\rho) :=
S(\rho\|\Delta[\rho])$, and for a Markovian semigroup $\rho_t =
e^{t\mathcal{L}}(\rho)$ the loss rate is $\dot{C}_{\mathrm{loss}}(\rho) :=
-\frac{d}{dt}C(\rho_t)|_{t=0}$.
\end{definition}

\subsection{Single-law lower bounds on maintenance power}

\begin{theorem}[Imported, corrected form; proved in \cite{MaintWork} (v2)]
\label{thm:imported}
Consider a finite-dimensional system at temperature $T$ undergoing uncontrolled
Markovian dynamics with thermal stationary state, controlled by stroboscopic
battery-assisted thermal operations. Then:
\begin{enumerate}
\item[(a)] (\emph{Unconditional}) The minimal total maintenance power obeys
$P_{\min}(\rho) \ge \kB T\,\sigma(\rho)$, where $\sigma(\rho) \ge 0$ is the
entropy production rate of the target, which splits as $\sigma =
\dot{F}^\Delta_{\mathrm{diag}} + \dot{C}_{\mathrm{loss}}$.
\item[(b)] (\emph{Coherence power}) Under diagonal contraction (automatic for
Davies generators), $P_{\min}(\rho) \ge \kB T\,\dot{C}_{\mathrm{loss}}(\rho)$.
\item[(c)] (\emph{Extra power}) Against thermodynamically
$\varepsilon$-efficient population-maintenance baselines, under population
autonomy, $P(s) - P(s_{\mathrm{diag}}) \ge \kB T\,\dot{C}_{\mathrm{loss}}(\rho)
- \varepsilon$; in the free-baseline case of \cite{MaintWork} --- e.g.\ pure
dephasing, where both $\mathcal{L}(\Delta[\rho]) = 0$ and
$\Delta[\mathcal{L}(\rho)] = 0$ --- one has $P_{\mathrm{extra}}(\rho) =
P_{\min}(\rho) \ge \kB T\,\dot{C}_{\mathrm{loss}}(\rho)$ with no further
assumptions.
\end{enumerate}
\end{theorem}

\begin{remark}[Correction relative to v1 of this note]\label{rem:importfix}
v1 imported the bound as: ``assume control strategies exist for $\rho$ and for
$\Delta[\rho]$; then $P_{\mathrm{extra}}(\rho) \ge \kB
T\,\dot{C}_{\mathrm{loss}}(\rho)$.'' In that generality the statement is
\emph{false}: with $P_{\mathrm{extra}}$ defined as an infimum over unlinked
strategy pairs (as in \cite{MaintWork} v1), the infimum is $-\infty$
(battery-burning baselines). The corrected hierarchy is
Theorem~\ref{thm:imported}(a)--(c); see \cite{MaintWork} (v2), whose v2 also
corrected the sign convention of work (battery consumption, not increase).
\end{remark}

\subsection{What these bounds do and do not require}

The bounds are purely operational: they depend on a control model and
dynamical functionals, not on geometric structure. They remain valid even if
all geometric scaling hypotheses fail. Geometry enters only when one asks how
$\dot{C}_{\mathrm{loss}}$ or effective rates depend on a separation parameter.

\section{Layer II: standard dynamical inputs}

In weak-coupling limits with secular approximation, generators take the Davies
form \cite{Davies74,BreuerPetruccione}; coherence sectors decouple and decay
at rates fixed by the generator; detailed balance yields the classical
H-theorem for populations (diagonal contraction) \cite{Spohn78}. These facts
support rate domination in controlled regimes, but do not by themselves imply
any dependence on geometric separation --- and, as \cite{Stress} (v2)
quantifies, the secular construction has a price: at finite size its jump
operators are delocalized, so ``local coupling'' statements within the Davies
class must be interpreted with care.

\section{Layer III: the hinge, updated from open to partially resolved}

\begin{hypothesis}[Rate inheritance principle (RIP), v1 formulation]
\label{hyp:rip}
Let $\kappa(\epsilon)$ denote an effective coherence-loss rate associated with
an operational separation parameter $\epsilon$ in a validated Markovian
regime. Rate inheritance holds if there exist $A, \beta > 0$ with
$\kappa(\epsilon) \le A\,\mathrm{poly}(m\epsilon)\, e^{-\beta m\epsilon}$
uniformly over the regime of interest.
\end{hypothesis}

\begin{remark}[Status, July 2026]\label{rem:status}
v1 identified Hypothesis~\ref{hyp:rip} as the unique unresolved hinge. Its
status is now:
\begin{itemize}
\item \textbf{Derived in an exactly solvable quasi-free local-sink class and
verified against exact Liouvillian spectra} \cite{HCut}: for a sub-gap
probe coupled through a gapped quasi-free buffer to a strictly local Markovian
sink, the exact Liouvillian rapidity obeys $\kappa(\epsilon) =
P(\epsilon)\,e^{-2q(\omega_b)\epsilon}$ with the frequency-resolved
Combes--Thomas exponent $\cosh q(\omega) = (\mu^2 + 4 - \omega^2)/(4\mu)$,
verified against exact spectra over ten decades; the exponent closes at the
band edge and vanishes in-band and at criticality.
\item \textbf{Failure scenario realized, Davies class} \cite{Stress}: with
near-zero-Bohr-frequency channels, the projected floor envelope persists under
separation --- v1's ``$\omega = 0$ floors'' scenario --- \emph{within} the
secular model class, whose jump operators are delocalized at finite size
($\sim 90\%$ of the zero-frequency component's weight beyond the coupling
site); the two results bracket the physics rather than contradicting each
other.
\item \textbf{Open}: the interacting gapped buffer with strictly local
dissipation --- the one version not implicit in known physics --- stated as an
explicit conjecture in \cite{HCut}.
\end{itemize}
\end{remark}

\begin{remark}[Proxy discipline (new in v2)]\label{rem:proxy}
v1 allowed $\kappa(\epsilon)$ to be ``any operationally extracted rate proxy
consistent with the Markovian description.'' The series has since learned, at
its own expense, that this is too permissive: the windowed proxy of
\cite{HCut} v1 measured arrival kinematics, not gap physics (its evidence was
withdrawn after a critical-point control), and the ceiling envelope of
\cite{Stress} v1 could not feed a maintenance no-go (a floor is required).
Valid extractions include Liouvillian rapidities, fixed-absolute-time
exponential fits with gap-removed controls and dissipation-free baselines, and
projected Dirichlet-form floors with declared operator subspaces
\cite{HCut,Stress}; the protocol-level codification of this discipline is
\cite{RIP} (v2).
\end{remark}

\begin{remark}[Falsification criterion, sharpened]
Hypothesis~\ref{hyp:rip} is falsified in a given regime if static correlations
decay exponentially while no envelope bound on a \emph{validly extracted}
$\kappa(\epsilon)$ (Remark~\ref{rem:proxy}) holds there. All claims of
geometric scaling of coherence-maintenance power depend on this hinge, now
with its validity domain partially charted.
\end{remark}

\section{Scenario analysis (exercised, not merely stated)}

\begin{itemize}
\item \textbf{RIP holds in a regime} (now actual for the quasi-free local-sink
class): geometric suppression of effective rates combines with
Theorem~\ref{thm:imported} to give collar-dependent power envelopes, e.g.\
$P_{\mathrm{extra}} \gtrsim \kB T\,K_\nu(m\epsilon)\,C^\Delta(\rho)$ under the
two-sided interface assumptions of \cite{MaintWork} (v2), and a logarithmic
isolation cost for coherence lifetime \cite{HCut}.
\item \textbf{RIP fails through $\omega = 0$ floors} (now realized within the
Davies class): maintenance remains costly at any separation; one obtains an
operational resource horizon rather than geometric protection
\cite{Stress}.
\item \textbf{Markovianity fails}: the Davies picture is invalid; the
maintenance bounds of Theorem~\ref{thm:imported} remain operational but the
dynamical model requires revalidation.
\end{itemize}

The design goal of v1 --- robustness under partial refutation --- has been
exercised twice in practice: the withdrawal of \cite{HCut} v1's numerical
evidence and the vacuity of \cite{MaintWork} v1's extra-power definition both
forced repairs, and in both cases the operational core (Layer I) survived
unchanged while the affected layer was rebuilt with stronger statements.

\section{What this program does not claim}

This framework does not propose modified dynamics, collapse mechanisms,
derivations of the Born rule, or ontological interpretations of the
quantum--classical transition. Any notion of a ``cut'' is operational, defined
solely by control feasibility under explicit resource constraints
\cite{HCut}.

\section{Conclusion}

We provided an architectural closure of the coherence-maintenance program,
separating proved operational power bounds from conditional
dynamical/geometric interfaces and open hypotheses. Coherence maintenance is a
well-defined operational resource with unavoidable dynamical cost, independent
of geometry. Rate inheritance --- v1's unique unresolved hinge --- is now
derived in an exactly solvable class and verified against exact Liouvillian
spectra, realized in its failure scenario within the Davies class, and open
precisely for interacting gapped buffers with local dissipation; the proxy-validation discipline that the series acquired along
the way is recorded here as part of the program's method.

\appendix

\section{Changes relative to version 1}

\begin{enumerate}
\item Theorem~\ref{thm:imported} restates the imported maintenance bound in
the corrected form of \cite{MaintWork} (v2); v1's import reproduced the
vacuous unlinked-pairs formulation (Remark~\ref{rem:importfix}).
\item The hinge status is updated from ``open'' to ``partially resolved''
(Remark~\ref{rem:status}), with the quasi-free proof and the Davies-class
floor realization; the missing case (interacting buffers) is delegated to the
conjecture of \cite{HCut}.
\item New proxy-discipline remark (Remark~\ref{rem:proxy}) distilling the
windowed-proxy withdrawal of \cite{HCut} and the ceiling/floor correction of
\cite{Stress}.
\item References updated to the v2 series with companion-repository paths;
\cite{Stress} added (absent from v1); the ``core technical reference''
author-profile pointer of v1 is replaced by the repository.
\item The RIP framing note \cite{RIP} (v2) is folded into the series map
(editorial box, Section~1, Remark~\ref{rem:proxy}): it states the weak form
of rate inheritance as a lemma under explicit hypotheses and codifies the
falsification protocol referenced here.
\end{enumerate}

\begin{thebibliography}{9}\small

\bibitem{Davies74} E.~B.~Davies, \emph{Markovian master equations}, Commun.
Math. Phys. \textbf{39} (1974), 91--110.

\bibitem{Spohn78} H.~Spohn, \emph{Entropy production for quantum dynamical
semigroups}, J. Math. Phys. \textbf{19} (1978), 1227--1230.

\bibitem{BreuerPetruccione} H.-P.~Breuer and F.~Petruccione, \emph{The Theory
of Open Quantum Systems}, Oxford University Press (2002).

\bibitem{HastingsKoma} M.~B.~Hastings and T.~Koma, \emph{Spectral gap and
exponential decay of correlations}, Commun. Math. Phys. \textbf{265} (2006),
781--804.

\bibitem{Clustering} L.~Eriksson, \emph{Clustering, Recovery, and Locality in
Algebraic Quantum Field Theory: Quantitative Bounds via Split Inclusions and
Modular Theory}, ai.viXra:2512.0060, v2 (2026; v1: 2025). Companion
repository, \texttt{papers/0060-clustering-recovery}.

\bibitem{MaintWork} L.~Eriksson, \emph{The Conditional Maintenance Work
Theorem: Operational Power Lower Bounds from Energy Pinching and a
Split-Inclusion Blueprint for Type III AQFT}, ai.viXra:2512.0061, v2 (2026;
v1: 2025). Companion repository, \texttt{papers/0061-maintenance-work}.

\bibitem{HCut} L.~Eriksson, \emph{The Heisenberg Cut as a Resource Boundary:
An Operational Outlook from Coherence Maintenance Costs}, ai.viXra:2512.0064,
v2 (2026; v1: 2025). Companion repository,
\texttt{papers/0064-heisenberg-cut}.

\bibitem{Stress} L.~Eriksson, \emph{Stress Testing the Rate Inheritance
Principle: Spectral Decoherence Rates and an Operational Resource Horizon},
ai.viXra:2512.0070, v2 (2026; v1: 2025). Companion repository,
\texttt{papers/0070-rate-inheritance}.

\bibitem{RIP} L.~Eriksson, \emph{The Rate Inheritance Principle: From Static
Correlations to Dynamical Decoherence Rates}, ai.viXra:2512.0072, v2 (2026;
v1: 2025). Companion repository, \texttt{papers/0072-rip-principle}.

\end{thebibliography}

\end{document}
